% Ampliación argumental. Insertar tras la introducción de excepcional.tex,
% antes de su primera subsección. Conserva íntegro el desarrollo existente.
\paragraph{Articulación de las coordenadas regionales y la incidencia de clausura.}
\label{inc-ampl-articulacion}
La función de esta prolongación se comprende a partir de una distinción
entre tres operaciones: conservar un registro, evaluar una coordenada y
extraer una relación de incidencia. Las tres actúan sobre información
producida por APP--TRIT--TPK. Su articulación hace explícita la tesis
estructural del artículo: una realización numérica posee una procedencia
operatoria, y determinados datos de esa procedencia admiten, además, una
realización combinatoria y reticular. Las ecuaciones permiten seguir ambos
recorridos con sus dominios y con la información que permanece en cada uno.

El punto de partida es el estado terminal con sus prolongaciones
compatibles. En él se conservan los residuos y cocientes de las dos hojas
APP, el régimen y la orientación TRIT, y el transporte TPK con su memoria.
Las regiones de clausura, propagación y autoescala se determinan en ese
dominio antes de que sus lectores establezcan las coordenadas
$\pi,e,\varphi$. Al mismo tiempo, el registro orientado de las doce
ventanas determina el vector firmado $U$. La aparición posterior de $K$
procede de la transformación integral demostrada en
\S\ref{subsec:k-hadamard}:
\begin{equation}
 U=\mathcal H_{12}K,\qquad
 K=\frac14\mathcal H_{12}U,\qquad
 \mathcal H_{12}^{2}=4I_{12}.
 \label{inc-ampl-registro}
\end{equation}
La congruencia que define $\mathcal L_H$ garantiza la integralidad de
la inversa. De este modo, $K$ constituye una coordenatización reversible
del registro firmado en la subred pertinente. Sus posiciones conservan el
orden de fase; sus diferencias $D_3K,D_4K$ conservan las relaciones
transversales; la suma $Q(K)$ determina la componente uniforme. Esa
estructura explica por qué el registro interviene tanto en la composición
numérica como en la selección de una configuración de frontera.

La evaluación periódica
\[
 K\longmapsto\kappa_{\mathrm{per}}
 =\frac{\sum_{m=1}^{12}K_m1000^{12-m}}{1000^{12}-1}
\]
mantiene recuperables los doce bloques cuando se fijan base, longitud y
origen de fase. La extracción de un soporte realiza una operación
diferente: conserva la pertenencia de cada posición a una clase definida
por un umbral y agrupa los registros que tienen el mismo soporte. La
distinción entre estas dos aplicaciones resulta decisiva. La precisión
escalar de $\kappa_{\mathrm{per}}$ y la incidencia de $B_K$ son dos
contenidos matemáticos específicos, relacionados mediante el registro
completo del que proceden.

\paragraph{La normalización de alfa y su configuración firmada.}
\label{inc-ampl-clausura}
La coordenada de $\alpha$ se determina en la vía posicional mediante la
composición de los bloques regionales y del registro dodecafásico. Antes
de la división posicional, la combinación orientada tiene componentes
$s_j=P_j+E_j-\Phi_j-K_j^{\mathrm{per}}$. La normalización satisface
\begin{equation}
 A_j=s_j+c_{j+1}-1000c_j,\qquad 0\leq A_j<1000.
 \label{inc-ampl-normalizacion}
\end{equation}
Esta identidad conserva simultáneamente el bloque normalizado y la
contribución del cociente transportado. El límite de los prefijos
compatibles determina la coordenada escalar de la clausura. La
configuración firmada conserva, por su parte, qué posiciones de la
combinación anterior a la normalización presentan defecto negativo.
En la ventana inicial, esa configuración es el vector $s_0$ de
\eqref{exc:precarry}, y su soporte negativo es $H^-_\alpha$.

Así quedan definidos dos contenidos de una misma operación. La
normalización posicional determina una coordenada real mediante sus
restos y cocientes compatibles; el soporte negativo determina una
configuración de posiciones en el ciclo de doce componentes. El signo se
refiere al balance de registros anterior a la normalización. Su lectura
incidencial permanece disponible porque el estado conserva ese balance
junto con el resultado normalizado. Esta articulación corresponde a la
vía posicional demostrada; la comparación con los operadores analíticos
de vacancias conserva el alcance preciso establecido en el capítulo de
$\alpha$.

\paragraph{De la codificación ternaria a la bandera marcada.}
\label{inc-ampl-bandera}
Las regiones aportan también los códigos $w_\pi,w_e,w_\varphi$ y
$w_{\rm APP}$ en el marco heredado. La aplicación
\[
 \gamma_{A_W}:w\longmapsto(w,wA_W)
\]
conserva la palabra de seis componentes como primera coordenada y añade
su transformación por $A_W$. La relación $A_W^2=-I_6$ fija la
compatibilidad cuadrática de esta realización. El paso siguiente,
$c\mapsto\operatorname{supp}(c)$, conserva las posiciones no nulas.
En las palabras de peso seis, el soporte identifica exactamente la pareja
$\{c,-c\}$, como prueba la distancia mínima del código. De esta forma,
el cociente por signo produce las hexadas, mientras el código completo
sigue disponible para las operaciones que requieren sus símbolos.

En esta carta, las dos construcciones de la cuaterna coinciden:
\begin{equation}
 \{i:K_i\geq729\}
 =H_e\cap P_\pi=B_K=\{4,6,9,10\}.
 \label{inc-ampl-cuaterna}
\end{equation}
El lado izquierdo procede de una comparación entera de los componentes
de $K$; el derecho, de la intersección de los soportes codificados de
propagación y clausura. El umbral pertenece a la carta de bloques
declarada y permanece explicitado en la definición del lector. El
registro íntegro conserva los valores sobre los que actúa, incluso cuando
distintos umbrales producen el mismo subconjunto. La igualdad expresa,
por tanto, una compatibilidad entre aplicaciones especificadas.

El soporte $H^-_\alpha$ completa esa cuaterna hasta una hexada. La
condición de pertenencia al soporte $P_\pi$, junto con los tamaños de
intersección $(4,3,2)$ respecto de $H_e,H_\varphi,H_{\rm APP}$,
selecciona una única hexada. Las pruebas siguientes reúnen las dos
determinaciones: la obtenida mediante el signo de $s_0$ y la obtenida
mediante incidencia. El resultado se organiza como
\begin{equation}
 \widetilde{\mathcal I}_*
 =\bigl(B_K\subset H^-_\alpha;
         P_\pi,H_e,H_\varphi,H_{\rm APP}\bigr).
 \label{inc-ampl-marco}
\end{equation}
La inclusión es la bandera; los cuatro soportes adicionales constituyen
su marco marcado. Las permutaciones simultáneas transportan el marco y
conservan sus inclusiones e intersecciones. La información relevante reside
en esta configuración transportada y en su clase de isomorfía.

Las cinco regiones de $\pi$ recuperan aquí su diferenciación interna.
Comparten $w_\pi$ como código ternario visible, pero conservan emisiones,
firmas y multiplicidades distintas. En la realización binaria sobre una
dodecada marcada, la región transversal corresponde a la octada cuya
intersección con la dodecada es la cuaterna; las otras cuatro regiones
corresponden a las elevaciones de las cuatro hexadas de su estrella. La
región negativa queda distinguida por $H^-_\alpha$ y el orden de carta
individualiza las tres regiones positivas. Esta correspondencia conserva
los papeles regionales y la distribución $432+4\cdot144=1008$, además
de realizar la partición incidencial $1+4$.

\paragraph{Del origen incidencial a la modificación reticular.}
\label{inc-ampl-reticulo}
La bandera marcada determina una posición mediante una operación
conjuntista equivariante:
\[
 o_*=(H_e\cap H_\varphi)
       \setminus(H^-_\alpha\cup H_{\rm APP})=\{1\}.
\]
Su función consiste en distinguir una componente entre las doce que
intervienen en la realización reticular. El código ternario completo
actúa como código de pegado en el módulo discriminante
$(A_2^*/A_2)^{12}$. Su preimagen define $N(C_W)$, y la
autodualidad, los pesos y la distancia del código se traducen,
respectivamente, en propiedades de integralidad y determinante, paridad
y cotas de norma de las clases de pegado. Aquí la información simbólica
del código cumple una función que el soporte aislado ya no desempeña.

La métrica $A_2$ y la cámara radial positiva se fijan en el desarrollo
reticular. Dentro de esa cámara, la congruencia necesaria para un vecino
par de índice tres tiene doce realizaciones de norma mínima $54$,
permutadas por el cambio de componente distinguida. El origen $o_*$
selecciona
\[
 v_\alpha=4\rho_1+\rho_2+\cdots+\rho_{12},\qquad
 v_\alpha^2=2(4^2+11)=54.
\]
El coeficiente $4$ pertenece a esa solución de la congruencia en la
cámara declarada. El subíndice $\alpha$ registra la procedencia
incidencial de la selección: su contenido es el marco que intervino en
la clausura y que ahora individualiza una dirección reticular.

El emparejamiento con $v_\alpha$ determina el subretículo
$N_{\alpha,0}$ de índice tres, y la adición de la clase
$v_\alpha/3$ produce el vecino $N_\alpha$ de
\eqref{exc:vecino}. Esta modificación conserva el rango y recupera
unimodularidad y paridad; la selección por congruencia elimina las raíces
anteriores, y las cotas locales excluyen raíces en las nuevas clases.
La identificación con Leech aparece después de esas verificaciones. La
realización binaria de las cinco regiones y esta construcción reticular
son dos prolongaciones del marco común: la segunda recibe directamente
el código ternario como dato de pegado.

\paragraph{Graduación, unidad y automorfismos.}
\label{inc-ampl-automorfismos}
La construcción de Frenkel--Lepowsky--Meurman se aplica al retículo
$N_\alpha$ una vez verificadas sus propiedades. Incorpora el espacio
de osciladores, la extensión central del retículo, el producto de vértice
y el sector torcido por el levantamiento de la negación. El resultado es
el álgebra $V^\natural_{\rm HMT}$ de \eqref{exc:flm}, cuyo grupo
de automorfismos es isomorfo al Monstruo. Su dominio de acción queda así
fijado por una estructura algebraica con graduación y producto. El teorema
de Moonshine de Borcherds determina posteriormente las trazas graduadas
con inserción de esos automorfismos \cite{flm,borcherds}.

La unidad adquiere en este nivel una realización precisa: la componente
de peso cero es la línea del vacío y aporta el coeficiente $1$ de
$q^{-1}$ en el carácter graduado. En el nivel cilíndrico, las
aplicaciones de refinamiento conservan la función constante mediante
precomposición. Cada afirmación de conservación tiene, por tanto, su
objeto y su aplicación. El desarrollo de las pruebas hace visible esa
continuidad de construcción, desde las operaciones aritméticas y la
memoria hasta la incidencia, la métrica y los productos de vértice.

Esta articulación sitúa a $K$ y a la clausura de $\alpha$ en el argumento
principal. $K$ conserva reversiblemente una coordenada integral de la
historia; la normalización de $\alpha$ compone los registros regionales
con esa coordenada; la configuración firmada determina una bandera; y
la bandera marcada selecciona la modificación reticular sobre la que se
realiza Moonshine. Las demostraciones siguientes desarrollan estas
aplicaciones, distinguiendo las equivalencias reversibles, los cocientes
de información y las realizaciones que añaden métrica o estructura
algebraica. Así se expresa una relación demostrable entre estructura,
ecuación y valor, conservando la atribución de cada construcción y el
alcance de sus teoremas.
